\documentclass[aspectratio=1610, 9pt]{beamer}
\setbeamertemplate{section in toc}[sections numbered]
\setbeamertemplate{subsection in toc}{\hspace{1em}\inserttocsectionnumber.\inserttocsubsectionnumber.~\inserttocsubsection\par}

\linespread{1.3}
\usetheme{default}

\usepackage{mathpazo}
\renewcommand{\familydefault}{\rmdefault}

\usepackage{graphicx}
\usepackage{amsmath}
\usepackage{amsthm}
\usepackage{amsfonts} 
\usepackage{mathtools}
\usepackage{bbm}
\usepackage{tikz}
\usepackage[normalem]{ulem}
\usepackage[dvipsnames]{xcolor}
\usepackage{setspace}
\usepackage{array}
\usepackage{booktabs}
\usepackage{caption}
\usepackage{subcaption}
\usepackage{dsfont}
\usepackage{natbib}
\usepackage{appendixnumberbeamer}

\usepackage{threeparttable}
\usepackage{booktabs}
\usepackage{multirow}
\newcommand{\sym}[1]{\ifmmode^{#1}\else\(^{#1}\)\fi}

\bibliographystyle{apalike}
\usepackage{hyperref}
\hypersetup{
    colorlinks=true,
    linkcolor=blue,
    urlcolor=blue,
    citecolor=blue,
}
\setbeamersize{text margin left=0.8cm, text margin right = 5mm}
\newtheorem{assumption}{Assumption}

% \setbeamertemplate{theorems}[numbered]
% \setbeamertemplate{footline}[frame number]

\newcommand{\indep}{\rotatebox[origin=c]{90}{$\models$}}
\newcommand\independent{\protect\mathpalette{\protect\independenT}{\perp}}
\def\independenT#1#2{\mathrel{\rlap{$#1#2$}\mkern2mu{#1#2}}}



\title[Chen 2025]{Anticipation Effects in Difference-in-Differences Design}

\author{Kyunghoon Ban, Zhengrun Chen, D\'esir\'e K\'edagni}




\begin{document}
\maketitle

\begin{frame}{Introduction}
\begin{itemize}
    \item The Difference-in-Differences (DID) relies on two main assumptions
         \begin{itemize}
     \item[-] Parallel Trends (PT)
     \item[-] No Anticipation Effects (No-AE), often imposed implicitly
     \end{itemize}
     
    \item People are forward-looking, e.g. 
     \begin{itemize}
     \item[-] Healthcare utilization decreases before Medicaid expansion (Alpert, 2016)
     \item[-] Physician labor supply increases before tort reform (Malani \& Reif, 2015)
     \end{itemize}
     \item Common practice: drop a few pre-treatment periods potentially affected by anticipation. 
     \begin{itemize}
         \item We don’t always know when people start to change behavior in anticipation of a policy.
         \item Different people/groups may react at different times.
     \end{itemize}
    \item Goal: Identify the ATT while allowing for violations of the No-AE assumption
\end{itemize}
\end{frame}

\begin{frame}{Identification}
\begin{itemize}
    \item Under the Parallel Trend, the standard DID estimand identifies the ATT net of the AE:
     $$\tau^{DID} = ATT-AE$$ where $AE\equiv\mathbb{E}[Y_0(1)-Y_0(0)\mid D=1]$
     \item Propose a bounded-variation assumption (Manski and Pepper, 2018; Roth and Rambachan, 2023): for some $c\in [0,1)$, $|AE| \leq c\cdot|ATT|$ 
     

     \item With multiple pre-treatment periods, we further assume exponential c-bounded anticipation: \\for some $c \in [0,1)$, 
        $|AE_{t_0}|\leq c^{|t_0|+1}\cdot |ATT|$ for all $t_0\in \{-T_0,...,-1,0\}$\\
        (Motivated by the exponential discounting model in Malani and Reif (2015))

        \item Interpretation:
    \begin{itemize}
        \item[-] Anticipation effects are bounded and decay over time
        \item[-] Provides a flexible, nonparametric alternative to fully parametric models
        \item[-] The choice of c can play the role of sensitivity analysis
    \end{itemize}
  
\end{itemize} 

\end{frame}


\begin{frame}{Summary}
\begin{itemize}
    \item We derive bounds on the ATT that take the form of intersection-union bounds
    \item We construct a misspecification robust bound for scenarios where researchers do not know whether the PT or No-AE assumption is the one failing.
    \item We extend the framework to staggered adoption design
    \item We illustrate the practical usefulness of these bounds through two applied empirical examples.
\end{itemize}


\end{frame}





\end{document}